% !TeX root = ./main.tex

\hfutsetup{
  title              = {OpenHarmony标准版适配启智01开发者套件（基于昇腾310B）并完成NPU推理功能},
  title*             = {Adapting OpenHarmony Standard to Qizhi01 Developer Kit (Based on Ascend 310B) with NPU Inference},
  author             = {潘孝成~~杨洋~~芦甲翔},
  author*            = {Pan Xiaocheng, Yang Yang, Lu Jiaxiang},
  stulist            = {2023213497~~2023213496~~2023213499},
  speciality         = {软件工程},
  supervisor         = {李琳~~张本宏},
  department         = {软件学院},
  % date               = {2025-07},
}

% 加载宏包

% 代码高亮
\usepackage{listings}
\usepackage{xcolor}

\renewcommand{\lstlistingname}{代码清单}
\definecolor{codebg}{RGB}{245,245,245}
\definecolor{codegreen}{RGB}{0,128,0}
\definecolor{codegray}{RGB}{128,128,128}
\definecolor{codeblue}{RGB}{0,0,200}
\definecolor{codepurple}{RGB}{128,0,128}

\lstset{
  backgroundcolor=\color{codebg},
  basicstyle=\ttfamily\small,
  breaklines=true,
  frame=single,
  framesep=5pt,
  rulecolor=\color{codegray},
  commentstyle=\color{codegreen},
  keywordstyle=\color{codeblue},
  stringstyle=\color{codepurple},
  numbers=left,
  numberstyle=\tiny\color{codegray},
  numbersep=8pt,
  tabsize=4,
  showstringspaces=false,
  captionpos=b,
  aboveskip=10pt,
  belowskip=10pt,
}

% Shell 命令样式
\lstdefinestyle{shell}{
  language=bash,
  basicstyle=\ttfamily\small,
  backgroundcolor=\color{codebg},
  frame=single,
  rulecolor=\color{codegray},
  commentstyle=\color{codegreen}\itshape,
  keywordstyle=\color{codeblue}\bfseries,
  stringstyle=\color{codepurple},
  showstringspaces=false,
  numbers=none,
  tabsize=4,
  breaklines=true,
}

% 定理类环境宏包
\usepackage{amsthm}

% 插图
\usepackage{graphicx}

% [H] 浮动体固定位置
\usepackage{float}

% 三线表
\usepackage{booktabs}

% 跨页表格
\usepackage{longtable}

% 算法
\usepackage[ruled,linesnumbered]{algorithm2e}

% SI 量和单位
\usepackage{siunitx}

% 参考文献使用 BibTeX + natbib 宏包
% \usepackage[sort]{natbib}
% \bibliographystyle{hfutthesis-numerical}

% 配置图片的默认目录
\graphicspath{{figures/}}

% 数学命令
\makeatletter
\newcommand\dif{%
  \mathop{}\!%
  \ifhfut@math@style@TeX
    d%
  \else
    \mathrm{d}%
  \fi
}
\makeatother
\newcommand\eu{{\symup{e}}}
\newcommand\iu{{\symup{i}}}

% xurl 允许 URL 在任意位置断行，避免长 URL 换行异常
\usepackage{xurl}

% hyperref 宏包在最后调用
\usepackage{hyperref}
\hypersetup{
  linkcolor  = black,
  citecolor  = black,
  urlcolor   = black,
}
